\chapter[Results and Discussions]{Results and Discussions}{Results and Discussions}\label{res}
This chapter presents the experimental evaluation of the proposed DAVS framework, consolidating the empirical findings reported in our research paper and aligning them with the end-to-end pipeline described in the earlier chapters. The goal of the evaluation is to demonstrate that communication efficiency, Byzantine robustness, and auditability can be achieved simultaneously for privacy-sensitive medical imaging federated learning.

\subsection{Experimental Setup (Summary)}
The system is evaluated on the PathMNIST medical imaging benchmark (MedMNIST v2) under non-IID client partitions across 10 simulated hospital nodes. A SimpleCNN model with 422,089 parameters is trained for 20 federated rounds with 2 local epochs per round. Unless stated otherwise, the malicious-client ratio is 20\% (2 out of 10 nodes) and attacks are applied from Round~1.

\subsection{Compared Configurations and Metrics}
We compare three configurations: (i) standard FL using FedAvg with no defense, (ii) FL augmented with a history-based trust mechanism (BlockDFL-style), and (iii) the full FL+DAVS pipeline (Random Projection sketching, amnesic DAVS committee selection, PBFT-inspired validation, and permissioned blockchain logging). Metrics include clean and poisoned-run accuracy, macro-F1, rounds-to-threshold convergence, attack success rate (ASR), committee purity (malicious inclusion rate), consensus success rate, and communication overhead.

\subsection{Baseline Convergence and Clean Accuracy}
In the benign setting, all configurations converge to comparable test accuracy (\(\sim\)83--85\%), indicating that DAVS and sketching impose negligible accuracy penalty. Under poisoning, however, the differences become significant: FL+DAVS maintains robust accuracy and avoids model collapse observed in undefended FL runs.

\begin{table}[h]
\centering
\caption{Overall Performance Comparison Across Three Configurations}
\begin{tabular}{|l|c|c|c|}
\hline
\textbf{Metric} & \textbf{FL} & \textbf{FL+BlockDFL} & \textbf{FL+DAVS} \\
\hline
Clean Test Accuracy (\%) & 83.22 & 82.05 & 82.41 \\
Macro-F1 Score (\%) & 83.90 & 82.63 & 82.74 \\
Rounds to 80\% Accuracy & 14 & 12 & 10 \\
Per-round Upload (MB) & 1.60 & 0.42 & 0.16 \\
Poisoned-Run Accuracy (\%) & 62.18 & 71.46 & 82.14 \\
Model Collapse in 10 Runs & 4 & 1 & 0 \\
\hline
\end{tabular}
\label{tab:primary_compare_ch8}
\end{table}

\subsection{Robustness Under Diverse Poisoning Attacks}
To evaluate Byzantine resilience, five poisoning strategies are tested: label-flip, gradient-flip, Gaussian noise injection, model scaling, and backdoor trigger insertion. Across all attack types, FL+DAVS achieves the highest robust accuracy and the lowest ASR (at most 11.8\%), while undefended FL exhibits large accuracy degradation.

\begin{table}[h]
\centering
\caption{Model Accuracy (\%) and ASR Under Each Attack Vector (20\% Malicious Clients)}
\begin{tabular}{|l|c|c|c|c|}
\hline
\textbf{Attack} & \textbf{FL} & \textbf{FL+BlockDFL} & \textbf{FL+DAVS} & \textbf{ASR (FL+DAVS)} \\
\hline
Label Flip & 64.80 & 73.55 & 82.47 & 8.6\% \\
Gradient Flip & 60.74 & 70.42 & 80.93 & 10.4\% \\
Gaussian Noise & 63.29 & 72.16 & 81.76 & 9.2\% \\
Model Scaling & 58.37 & 68.91 & 79.52 & 11.8\% \\
Backdoor Trigger & 66.02 & 74.63 & 83.12 & 6.9\% \\
\hline
\end{tabular}
\label{tab:attack_20_ch8}
\end{table}

\subsection{Committee Selection Quality and Consensus Reliability}
The key advantage of DAVS is its ability to form a clean validator committee using only current-round statistics. Under 20\% malicious participation, history-based selection admits malicious nodes into the committee in multiple rounds, whereas DAVS achieves 0\% malicious inclusion across all 20 rounds. This directly improves the reliability of PBFT validation and eliminates consensus failures in our experiments.

\begin{table}[h]
\centering
\caption{Committee Quality Comparison Over 20 Rounds (20\% Malicious)}
\begin{tabular}{|l|c|c|}
\hline
\textbf{Selection Metric} & \textbf{FL+BlockDFL} & \textbf{FL+DAVS} \\
\hline
Avg.\ malicious inclusion (\%) & 13.6 & 0.0 \\
Rounds with 0 malicious selected & 7/20 & 20/20 \\
Honest-malicious score gap & 0.29 & 0.68 \\
Consensus failures / 20 rounds & 3 & 0 \\
\hline
\end{tabular}
\label{tab:committee_quality_ch8}
\end{table}

\subsection{Communication and System Efficiency}
Random Projection sketching reduces the 422,089-dimensional updates to 128-dimensional sketches (approximately \(3{,}297{:}1\) compression), decreasing per-round upload from 1.60 MB to 0.16 MB. Committee-scoped PBFT reduces message complexity by restricting consensus to \(k{=}5\) validators. Overall, the framework reduces total bandwidth while maintaining a 100\% consensus success rate.

\begin{table}[h]
\centering
\caption{System-Level Efficiency Metrics}
\begin{tabular}{|l|c|c|c|}
\hline
\textbf{Metric} & \textbf{FL} & \textbf{FL+BlockDFL} & \textbf{FL+DAVS} \\
\hline
Consensus success rate & N/A & 96.5\% & 100.0\% \\
Avg.\ round latency (s) & 1.42 & 1.67 & 1.72 \\
Total bytes (20 rounds, MB) & 32.0 & 8.4 & 3.2 \\
Poisoned-run collapses & 4/10 & 1/10 & 0/10 \\
Auditability & No & Yes & Yes \\
\hline
\end{tabular}
\label{tab:efficiency_ch8}
\end{table}

\subsection{Discussion and Key Findings}
The experiments establish three core outcomes: (i) amnesic, round-local trust evaluation (cosine similarity + magnitude deviation) is more reliable than history-based reputation under adaptive adversaries, (ii) the full DAVS pipeline sustains robust accuracy within 3--4\% of the clean baseline across diverse poisoning attacks and up to 30\% malicious clients, and (iii) these security gains are achieved alongside large communication savings enabled by sketching and committee-scoped consensus, while providing immutable provenance through blockchain logging.

\section{Outputs}
\begin{figure}[h!]
    \centering
    \includegraphics[width=0.8\textwidth]{Chapter1/dashboard.png} % Adjust the width as needed
    \caption{Baseline FedAvg Training Output (Representative Logs/Curve)}
    \label{f1} % Optional: for referencing the figure in your document
\end{figure}

\begin{figure}[h!]
    \centering
    \includegraphics[width=0.8\textwidth]{Chapter1/ss2.png} % Adjust the width as needed
    \caption{Gradient Sketching Compression Summary (Representative Output)}
    \label{f1} % Optional: for referencing the figure in your document
\end{figure}

\begin{figure}[h!]
    \centering
    \includegraphics[width=0.8\textwidth]{Chapter1/ss4.jpeg} % Adjust the width as needed
    \caption{DAVS Committee Selection Behaviour (Representative Output)}
    \label{f1} % Optional: for referencing the figure in your document
\end{figure}

\begin{figure}[h!]
    \centering
    \includegraphics[width=0.8\textwidth]{Chapter1/ss5.png} % Adjust the width as needed
    \caption{PBFT Validation Phase Output (Representative Logs)}
    \label{f1} % Optional: for referencing the figure in your document
\end{figure}

\begin{figure}[h!]
    \centering
    \includegraphics[width=0.8\textwidth]{Chapter1/ss6.png} % Adjust the width as needed
    \caption{Blockchain Audit Trail Integrity Verification (Representative Output)}
    \label{f1} % Optional: for referencing the figure in your document
\end{figure}


%\chapter{Conclusions \& Future Scope}
